%%%%%%%%%%%%%%%%%%%%%%%%%%%%%%%%%
%
%       EXERCÍCIO
%
%%%%%%%%%%%%%%%%%%%%%%%%%%%%%%%%%

\ifdefstring{\atividade}{prova}{%
    \ifdefstring{\modo}{objetivo}{%
        \renewcommand{\valorquestao}{\ValorQObj\ ponto}
    }{%
        \renewcommand{\valorquestao}{\ValorQDisc\ pontos}
    }%
}%

\begin{exercicioBanco}[\valorquestao]
Seja \(f: \bbR \to \bbR\) uma função tal que \(f(\frac{1}{3}) = \sqrt{a}\), com \(a\) um número real positivo, e \(f(x+1) = x f(x)\), para todo \(x \in \bbR\). Determine \(f(\frac{7}{3})\).

% Define as alternativas
\newcommand{\alternativas}{%
    \begin{center}
        \begin{tabularx}{\textwidth}{XXXXX}
            (a) \(\dfrac{4\sqrt{a}}{9}\). &
            (b) \(\dfrac{10\sqrt{a}}{9}\). &
            (c) \(\dfrac{28\sqrt{a}}{9}\). &
            (d) \(\dfrac{40\sqrt{a}}{9}\). &
            (e) \(\dfrac{70\sqrt{a}}{9}\).
        \end{tabularx}
    \end{center}
}

% Define a resposta correta
\newcommand{\resposta}{A}

% Lógica condicional para exibição
\ifdefstring{\atividade}{lista}{%
    \alternativas
    \vspace{0.5em}
    
    \noindent\textbf{Resposta:} letra \textbf{\resposta}.
}{%
    \ifdefstring{\modo}{objetivo}{%
        \alternativas
    }{%
        % modo = discursiva → não mostra alternativas
    }
}
\end{exercicioBanco}

